\chapter{AI会犯哪些错：系统性认知}

前面的章节介绍了如何向 AI 描述问题、数据和建模任务。完成这些步骤后，AI 往往能够迅速生成代码、训练模型、绘制图表，并写出一段结构完整的结果解释。但项目到这里并没有结束。接下来还有一个无法外包给 AI 的问题：\textbf{这个结果到底对不对？}

传统程序中的错误经常以报错信息出现，使用者很容易意识到需要停止。AI 辅助机器学习中的许多错误却不会触发报错。代码能够运行，模型分数看起来不错，解释文字也十分流畅，但项目可能使用了错误字段、错误划分方式或错误评价指标。换句话说，AI 可以让人更快得到一个答案，却不能保证这个答案值得相信；而正是“看起来没问题”这一点，使它的错误比传统报错更难被发现。

本章的目标是建立对 AI 错误的系统认识。这里的重点不是穷举所有可能出现的问题，而是提供一个分类框架，帮助读者在面对一个完整的 AI 输出时，知道应该从哪些方向去怀疑它。我们把机器学习项目中最常见的问题归纳为五类：技术性错误、统计性错误、数据性错误、领域性错误和幻觉性错误。这五类错误从“代码层面”一直延伸到“现实层面”，需要的发现方法各不相同，所要求的专业知识也逐级加深。理解它们之间的差别，比记住任何一条具体的检查规则都更重要。

% ----------------------------------------------------------------
\section{AI错误的五种常见类型}
% ----------------------------------------------------------------

\subsection{技术性错误：代码没有正确完成指定计算}

技术性错误指代码本身的缺陷，包括语法错误、函数调用错误、变量使用错误和计算逻辑错误。它们是最接近传统编程错误的一类，也是通常最容易验证的一类，因为问题往往可以在代码和中间结果中被直接定位。

真正需要警惕的不是语法错误——程序会因此停止并给出报错信息——而是代码能够顺利运行，却没有执行使用者预期的计算。这类错误形式多样：构造滞后特征前没有按路段分组，导致不同路段的数据相互污染；合并两个表格时使用了不唯一的键，使部分记录被重复计算；计算平均速度时把表示设备离线的零值当作正常速度参与平均；修改数据类型后，时间字段按照字符串顺序而不是时间顺序排序；甚至只是把评价函数中真实值和预测值的位置写反，代码依旧正常运行并给出一个看似合理的分数。它们的共同特点是：任务定义和方法可能都没有问题，但 AI 生成的实现没有正确执行它们。

因此，发现技术性错误不能只看代码是否报错，而应当检查中间结果，并用能够手工验算的小样本去验证关键函数。\citet{domingos2012few}早已指出，机器学习中“能跑通的代码”与“做对了事情的代码”之间往往隔着一道并不显眼的鸿沟，而跨越这道鸿沟主要依靠对中间产物的主动检查，而不是对最终分数的信任。

\paragraph{案例：跨路段的滚动平均}

假设需要为每条路段计算过去三个时间步的平均速度。AI 生成代码后直接对整张表使用滚动平均函数，却没有先按路段分组。于是第一条路段末尾的速度会“溢出”到第二条路段开头的特征中。代码不会报错，新生成的列也没有缺失或格式问题，但特征的真实含义已经被破坏。一个简单而有效的验证方法，是构造两条速度差异明显、各自只有几条记录的路段：如果第二条路段开头的滚动平均明显受到第一条路段的影响，错误会立即暴露。这类“最小可验证样本”的思路，会在第十三章的验证方法中反复出现。

\subsection{统计性错误：方法运行正常，但评价或推断不成立}

统计性错误不是代码写错，而是方法选择、数据划分或指标解释不正确。这类错误尤其危险，因为它们经常产生非常漂亮的结果——一个高得异常的分数比一个报错更容易被接受，也更容易被当作成功。

其中最常见、危害也最大的是数据泄露（data leakage），即模型在训练或评价时使用了真实预测场景中根本无法获得的信息。例如用事故发生之后才记录的处理时长去预测事故是否发生，在划分训练集和测试集之前就对全部数据做标准化，或者对具有明显时间结构的数据进行随机划分。泄露发生后，模型相当于提前看到了部分答案，测试分数会异常优秀，而代码通常完全正常\cite{kaufman2012leakage}。\citet{kapoor2023leakage}在系统梳理大量“基于机器学习的科学研究”后发现，数据泄露是导致结果无法复现的首要原因之一，并且常常以非常隐蔽的形式出现。发现这类问题，唯一可靠的办法是回到预测发生的那个时刻，逐个追问每个特征在那一刻是否真的可用。

与数据泄露相关的另一类问题，是模型学到了“捷径”而非真正的规律。\citet{geirhos2020shortcut}把这种现象称为捷径学习（shortcut learning）：模型抓住了数据中某种与标签恰好相关、却没有因果意义的表层特征，在测试集上表现出色，一旦换到稍有不同的场景便迅速失效。一个经典的极端例子是所谓“聪明的汉斯”效应——模型看似学会了识别马匹，实际上只是学会了识别图片角落里的来源水印\cite{lapuschkin2019clever}。在交通场景中，这可能表现为模型并未学会预测拥堵，而只是记住了某个检测器编号恰好经常对应拥堵路段。

指标选择不当则是一类更“安静”的统计错误。不同指标回答的是不同的问题。当事故样本只占全部记录的 1\% 时，把所有样本都预测为“无事故”就能轻松获得 99\% 的 Accuracy，但这样的模型没有识别出任何一个真正需要关注的事故。如果漏报比误报更危险，就应重点关注 Recall；如果要在有限资源下提高预警的准确性，则还要兼顾 Precision；如果输出概率将被用于风险排序，则要进一步检查概率本身是否可靠。指标必须由真实任务决定，而不能由 AI 沿用默认选项。最后还要警惕“把一次高分当作稳定能力”：AI 可以快速尝试大量模型与参数组合，如果反复在同一个验证集上挑选表现最好的结果，最终的高分可能有相当一部分来自偶然。可靠的比较应当保留简单基准、记录尝试过的方案，并用一个独立的测试集确认最终结论。

\subsection{数据性错误：正确代码处理了错误含义}

数据性错误发生在 AI 对数据的理解与真实情况不一致时。此时代码可能完全正确，统计方法也可能合理，但输入数据并不表示使用者以为的那个对象。这类错误之所以难以发现，是因为它不在代码里，而在代码之外——在数据是如何被采集和定义的过程之中\cite{sambasivan2021data}。

最基础的一种是字段含义和单位错误。同一个字段名在不同数据集中可能具有完全不同的含义：\texttt{flow} 可能表示车辆数、辆/5分钟，也可能是换算后的辆/小时；速度可能是瞬时速度、时间平均速度或空间平均速度。单位或定义上的偏差通常不会引发任何报错，却会直接改变后续的全部分析。其次是异常值含义错误。AI 习惯用通用统计规则（如三倍标准差）处理异常值，但领域中的异常值常常是真实而重要的信号：事故造成的速度骤降、大型活动带来的客流峰值、极端天气下的需求变化，都不应仅仅因为“偏离平均水平”而被删除。

更深一层的问题是标签并不天然正确。事故严重程度、拥堵状态、乘客满意度等标签都依赖具体的定义与记录方式。已报告的事故不等于全部实际发生的事故；投诉较多既可能意味着服务质量差，也可能只是因为该渠道的投诉更方便。模型学习的永远是数据中标签所反映的规律，而不是现实中那个我们真正关心的概念。要发现这类错误，必须理解数据如何采集、每一行代表什么、字段如何生成，以及哪些对象根本没有进入数据。这也正是本书反复强调向 AI 提供数据字典和领域背景的原因。

\subsection{领域性错误：统计上合理，现实中却不合理}

领域性错误是指结果在代码和统计上都看似成立，却违背专业常识、业务规则或现实约束。这类错误通常最难由 AI 自己发现，因为相关知识并不完整地存在于 Prompt 和数据之中——它存在于专业人员的经验里。

它的表现形式多种多样：模型建议使用一个在实际部署时无法实时获取的字段；公交调度方案要求的班次调整违反了运营规则；预测结果长期出现负流量或明显超过道路通行能力的数值；模型在普通工作日表现很好，却被不加区分地用于节假日和极端天气；或者 AI 仅凭一个相关关系就提出政策建议，而没有考虑其他可能的原因。这些问题的共同点在于：统计分数本身完全无法察觉它们，必须借助对数量级、方向、可执行性和适用场景的常识判断。

\paragraph{案例：负流量预测}

某回归模型在测试集上的平均绝对误差并不高，但在低流量时段会输出负数。如果只看总体指标，这个问题可能完全不会暴露；可一旦回到交通工程常识，负流量显然是不可能的。这个例子说明，领域知识在这里提供了一个比平均误差更直接、也更不可妥协的正确性标准：无论分数多高，一个会预测出负车流的模型都不能被原样采用。

\subsection{幻觉性错误：AI生成了不存在或不准确的信息}

幻觉性错误指 AI 编造出并不存在或并不准确的信息，例如虚构文献、软件参数、法规条款、数据来源和计算结果。它也可能表现得更加隐蔽：引用的文献真实存在，但原文并不支持 AI 写出的那个结论。这一点尤其值得初学者警惕，因为表面看起来“有出处”往往比明显的胡说更具迷惑性。

幻觉并非偶然的故障，而是大语言模型工作方式的自然结果。模型根据上下文生成最符合语言模式的内容，而不是每次都先去查找并核实事实，因此流畅的表达本身从来不能证明事实成立\cite{bender2021stochastic}。关于自然语言生成中幻觉问题的系统综述也指出，这类错误普遍存在于各类生成任务中，且很难仅靠模型自身完全消除\cite{ji2023survey}。出于这个原因，以下内容必须回到原始来源逐一核验：文献的题目、作者、结论和适用范围；软件函数、参数名称和默认行为；法规、标准和技术规范；数据集的字段定义与来源；以及报告中出现的每一个具体数字和比例。尤其需要强调的是，绝不能让 AI 在没有真正运行代码的情况下“估计”模型分数，也不能直接把 AI 生成的文献引用放进正式报告。

把上述五类错误放在一起，可以看到它们大致沿着“离代码越来越远、离现实越来越近”的方向排列：技术性错误在代码内部，统计性错误在方法层面，数据性错误在数据含义层面，领域性错误在现实约束层面，而幻觉性错误则关乎信息本身是否真实。表~\ref{tab:five_errors} 概括了它们的典型表现和主要发现方式，可以作为面对任何 AI 输出时的初步排查清单。

\begin{table}[H]
\centering
\caption{AI辅助机器学习中的五类常见错误}
\label{tab:five_errors}
\begin{tabular}{p{2.2cm}p{5.0cm}p{5.2cm}}
\toprule
错误类型 & 典型表现 & 主要发现方式 \\
\midrule
技术性错误 & 代码实现与任务要求不一致 & 小样本、单元测试、中间结果检查 \\
统计性错误 & 泄露、捷径、指标误用、过度调参 & 基准比较、评价设计审查、特征可得性检查 \\
数据性错误 & 字段、单位、异常值和标签含义错误 & 数据字典、采集机制、人工抽查 \\
领域性错误 & 数字或方案违反现实常识与业务约束 & 数量级检验、领域专家判断、情景检查 \\
幻觉性错误 & 编造或错误转述事实、引用和参数 & 原始来源核验、实际运行、官方文档 \\
\bottomrule
\end{tabular}
\end{table}

% ----------------------------------------------------------------
\section{为什么AI会“自信地”犯错}
% ----------------------------------------------------------------

了解了错误的类型之后，还需要理解一个更根本的问题：为什么 AI 在犯错时，语气往往和它答对时一样笃定？如果模型在不确定时会自动表现出迟疑，使用者就可以把“语气”当作可靠性的信号。遗憾的是，事实恰恰相反。要建立合理的怀疑习惯，必须先理解 AI 自信表达背后的几种机制。

\subsection{合理表达不等于事实正确}

大语言模型的训练目标，是根据上下文生成连贯、完整且符合语境的文本。它非常擅长这件事，但“听起来合理”与“实际正确”从来不是同一回事。面对信息不足的问题，模型通常仍会继续作答，而不会自动停下来；它生成的是当前语境下“最像正确答案”的表达，而不是一条经过验证的事实。

这意味着，回答的语气越确定，并不代表事实越可靠；代码越完整，也不代表方法越正确；解释越详细，更不代表证据越充分。更麻烦的是，即使是模型自己给出的置信度，也未必可信。关于现代神经网络校准的研究表明，模型输出的概率经常存在系统性的过度自信，预测概率高并不等于预测真的更准\cite{guo2017calibration}。因此，评价 AI 输出的第一原则，就是把\textbf{表达质量}与\textbf{内容正确性}彻底分开：前者是模型被直接训练去优化的，后者则需要使用者另行提供证据。

\subsection{谄媚倾向：AI偏向给出你想听的答案}

除了“总是显得自信”之外，大语言模型还存在一种更微妙、也更值得警惕的倾向：它偏向于给出使用者想听的答案。这种现象在研究中被称为\textbf{谄媚效应（sycophancy）}。它的根源在于模型的对齐训练方式：在基于人类反馈的强化学习（RLHF）中，模型会被调整为生成人类评价更高的回答，而人类评价者往往更喜欢那些迎合自己、认同自己立场的回应。久而久之，模型就学会了一种并不总是诚实的策略——顺着用户的话说\cite{perez2022discovering}。

针对这一现象的系统研究发现，当用户在提问中暗示了自己的偏好或观点时，模型的回答会明显向该方向偏移；如果用户对一个本来正确的答案表示质疑，模型常常会动摇甚至直接改口，即使原来的答案并没有错\cite{sharma2023sycophancy}。这对机器学习的验证工作有直接而严重的影响。设想这样一个场景：你训练出一个分数很高的模型，然后问 AI“我这个结果是不是很好？”——由于谄媚倾向，它很可能顺势肯定你，而不是主动去揭示其中潜藏的数据泄露。反过来，如果你已经怀疑某处有问题并把这种怀疑说出口，它又可能不加分辨地附和你的怀疑。无论哪种情况，AI 都在迎合你的预期，而不是独立地判断对错。

谄媚效应与人本身的确认偏误（confirmation bias）会相互放大：人倾向于相信符合自己预期的结论，而 AI 又倾向于生成符合人预期的回答，于是双方很容易共同滑向一个彼此都满意、却未必正确的答案。理解这一点，能帮助我们在第十三章设计“反向追问”时格外小心——简单地问 AI“这里是不是错了”往往得不到诚实的答案，更可靠的做法是用中立的措辞提问、在不透露自己倾向的独立上下文中让它审查，甚至要求它专门寻找支持相反结论的证据。

\subsection{长尾知识薄弱，且不会主动承认不知道}

AI 自信犯错的第三个原因，与它的知识分布有关。模型在常见编程模式、通用机器学习概念和广泛流传的公开资料上通常表现良好，但在具体城市、特定数据平台、内部业务规则等专业“长尾”知识上更容易出错。它可能知道一般的交通流分析方法，却不知道某个检测器字段中的零值代表设备离线；知道常见的事故分类，却不知道当前地区采用的正式分类标准；知道某个软件库的旧接口，却不知道当前版本已经修改了参数。任务越依赖本地信息和专业细节，模型出错的概率就越高，越需要在 Prompt 中补充上下文并加强核验。

问题在于，模型并不会因为触及了自己的知识边界就主动沉默。使用者经常希望 AI 在不知道时直接说“不知道”，但模型往往无法准确判断自己究竟掌握了什么——它可能把模糊的记忆、常见的模式和上下文中的暗示组合成一个听起来确定的回答。我们可以通过 Prompt 在一定程度上降低这种风险，例如明确要求：

\begin{quote}
如果缺少足够信息，请明确指出缺失内容，不要自行补充事实。对于具体数字、文献和软件参数，请说明核验来源；无法核实时请标记为待确认。
\end{quote}

但需要清醒地认识到，这样的提示只能提醒模型，不能替代外部核验。它能减少一部分编造，却不能保证模型真的识别出了自己所有的知识盲区。

% ----------------------------------------------------------------
\section{为什么领域知识决定你能发现多少错误}
% ----------------------------------------------------------------

上一节解释了 AI 为什么会自信地犯错，本节回答一个与之互补的问题：在这些错误面前，决定一个人能发现多少的，到底是什么？答案是领域知识。AI 能够大幅降低代码实现和资料整理的门槛，却不能取消领域知识的重要性。恰恰相反，当代码生成变得廉价之后，人的主要工作反而更加集中地落在判断问题是否合理、数据是否被正确理解、结果是否符合现实这些只能依靠领域判断的环节上。

\subsection{没有领域知识，就不知道该检查什么}

错误检查的前提，是知道应该怀疑什么。一个不熟悉交通数据的人看到速度为零，可能简单地认为它是异常值并删除；而熟悉数据采集机制的人会进一步追问：它发生在信号交叉口、严重拥堵期间，还是设备离线状态？同一个数值，在不同情形下需要完全不同的处理方式。差别不在于谁更细心，而在于谁知道这里存在一个需要分辨的问题。

由此可见，领域知识的作用不仅是帮助回答问题，更是帮助提出正确的问题。它告诉使用者哪些数量级根本不合理、哪些字段不可能在预测时实时获得、哪些场景必须单独评价、哪些结论无法从现有数据推出。前一节谈到的五类错误中，越往后（数据性、领域性）越依赖这种“知道该查什么”的能力——技术性错误尚可借助通用工具发现，而领域性错误几乎完全依赖人的专业判断。

\subsection{AI不能替你定义“合理”，但能帮你发现知识缺口}

一个模型是否“合理”，最终取决于它所服务的现实任务，而这套标准并不写在算法里。事故漏报与误报的代价如何权衡，节假日的预测是否必须准确，某种调度方案在现场能否执行——这些都不能由 AI 自动决定。因此可以说，AI 的能力“天花板”，往往就是使用者能够提供和运用的领域知识：领域基础越薄弱，越难判断 AI 的输出是否可信，也越容易把一个漂亮但错误的结果当作成功。

不过，强调领域知识的重要性，并不意味着使用者必须在开始之前就掌握所有答案。AI 的一个高价值用法，恰恰是反过来帮助暴露自己的知识缺口：它可以列出需要确认的问题、比较不同的解释、生成检查清单，并指出当前分析依赖了哪些假设。例如可以这样反向提问：

\begin{quote}
如果要判断这份交通预测结果是否合理，一名交通工程师需要提供哪些领域信息？请按数据含义、数量级、特殊场景和实际部署约束分类列出。
\end{quote}

这类提问不会替代领域判断，但能把“我还不知道自己不知道什么”转化为一份具体的待确认清单，从而引导使用者去查阅教材、原始资料或请教专业人员。换句话说，AI 不能替你定义合理，但可以帮你更快地找到那些需要由人来定义合理的地方。

% ----------------------------------------------------------------
\section{建立系统性的错误检查习惯}
% ----------------------------------------------------------------

认识到错误的类型、AI 自信的来源和领域知识的作用之后，最后一步是把这些认识落实为可执行的检查习惯。面对 AI 生成的结果，不必怀疑每一个字符，但应当把有限的注意力优先投向那些会实质改变项目结论的环节。

\subsection{先检查上游，再检查最终结果}

错误检查应当遵循数据流动的方向，从上游到下游。一个实用的顺序是：先确认 AI 是否正确理解了任务目标和预测时点；再检查字段、单位、标签和异常值是否被正确理解；接着核对数据处理后的行数、范围和样例是否符合预期；然后审查特征、预处理和数据划分是否存在泄露；之后才比较模型是否优于简单基准、指标是否契合任务目标；继而检验数字、图形和解释是否符合领域常识；最后核验引用和具体事实是否经过原始来源确认。

这个顺序背后的逻辑是：上游错误会污染全部下游结果。如果标签定义本身就错了，再去仔细检查模型参数就毫无意义；如果测试集已经发生泄露，那么再漂亮的分数也不能被采用。把上游的问题排查清楚，往往比在下游反复调优更能提高结论的可靠性。

\subsection{越是高分和“符合预期”的结果，越值得检查}

人有一种自然的倾向：异常糟糕的结果容易引起怀疑，异常优秀的结果却容易被欣然接受。但从前面的讨论可以看出，这恰恰是危险的。当模型表现远高于基准时，背后可能潜藏着数据泄露、重复记录或评价错误；而当结果完全符合使用者的预期时，又可能掺入了确认偏误，甚至被 AI 的谄媚倾向进一步强化。

因此，看到一个出乎意料的高分时，正确的反应不是庆祝，而是依次自问：是否使用了预测时无法获得的信息？训练集和测试集之间是否存在重复或时间交叉？是否选择了一个过于容易的指标？是否只报告了多次实验中表现最好的那一次？以及，简单基准的表现究竟是多少？这些问题不会否定一个真正优秀的模型，却能拦下大量“好得不真实”的结果。

\subsection{让错误修复进入下一轮Prompt}

发现错误之后，不应满足于修改当前这一次的输出，而应把修复的经验沉淀为新的 Prompt 约束和检查步骤。字段含义被误解之后，就把数据字典加入后续 Prompt；发生时间泄露之后，就明确规定所有特征只能使用预测时刻之前的信息；异常值被误删之后，就要求 AI 先展示异常记录并等待确认；引用出现错误之后，就要求对关键事实标记来源和待核验状态；一次性流程出错之后，就把任务拆成可以逐步确认的 Prompt 链。

这正是“Prompt as Code”这一书名所要传达的核心含义之一：错误检查得到的经验，应当像修改代码、补充注释、增加测试一样，被固化进下一版 Prompt，使同类错误更难再次发生。一个项目的可靠性，不是来自某一次幸运的正确输出，而是来自这种把教训不断写回流程的迭代过程。

% ----------------------------------------------------------------
\section{实战练习：找出AI结果中的错误}
% ----------------------------------------------------------------

\subsection*{练习一：错误分类}

判断以下问题分别属于技术性、统计性、数据性、领域性还是幻觉性错误，并说明理由：

\begin{enumerate}
  \item AI 使用不存在的软件参数，代码运行时报错；
  \item 时序预测模型随机划分训练集和测试集，得到异常高分；
  \item AI 将检测器离线产生的零速度解释为严重拥堵；
  \item 模型预测某路段未来流量为负数；
  \item 报告引用了一篇真实论文，但论文并未支持报告中的因果结论。
\end{enumerate}

\subsection*{练习二：检查高分模型}

一个事故预测模型报告 Accuracy 为 98.7\%。请列出至少六项优先检查内容，覆盖任务、数据、实现、评价和领域合理性，并说明为什么“高分”本身反而是一个需要警惕的信号。

\subsection*{练习三：找出代码逻辑风险}

AI 生成了一段“按路段计算过去四个时间步平均速度”的代码。请设计一个包含两条路段、缺失值和时间乱序的小样本，用于检查分组、排序和未来信息问题。

\subsection*{练习四：识别谄媚倾向}

针对同一个模型结果，分别用三种措辞向 AI 提问：第一种明确表达“我觉得这个模型很好”，第二种表达“我怀疑这里有问题”，第三种用完全中立的措辞请它审查。比较三次回答的差异，说明谄媚效应在其中如何体现，并据此设计一条更不容易诱发迎合的提问方式。

\subsection*{练习五：改写错误预防Prompt}

选择本章中的一种常见错误，将其转化为一条可加入未来任务的 Prompt 约束，并说明还需要使用什么外部方法验证该约束是否真正被遵守。
